\section{Conceptos fundamentales}

\textbf{Muestra:} Parte representativa de un material que se va a analizar. 

\textbf{Muestreo:} Proceso de selección y obtención de una muestra representativa de un todo.

\textbf{Analito:} Componente de la muestra que se desea identificar o cuantificar.


\section{Tipos de muestreo}

\subsection{Según su estado físico}

\begin{itemize}
    \item \textbf{Sólido:} Requiere técnicas como cuarteo, molienda o toma de núcleos.
    \item \textbf{Líquido:} Se utilizan métodos como el muestreo simple, compuesto o por bombeo.
    \item \textbf{Gaseoso:} Se usan trampas, adsorbentes o recipientes para la captura.
\end{itemize}

\subsection{Tipos de muestra}

\begin{itemize}
    \item \textbf{Muestra simple:} Tomada de una sola vez en un punto específico.
    \item \textbf{Muestra compuesta:} Mezcla de varias muestras simples, tomadas en diferentes puntos o tiempos.
\end{itemize}

\section{Plan de muestreo}

El plan de muestreo es un documento que detalla cómo, dónde y cuándo se tomarán las muestras. Debe incluir:

\begin{itemize}
    \item \textbf{Objetivo del muestreo:} ¿Qué se quiere saber?
    \item \textbf{Puntos de muestreo:} Ubicación geográfica o dentro de un proceso.
    \item \textbf{Número de muestras:} Determinado por métodos estadísticos para asegurar la representatividad.
    \item \textbf{Equipos y contenedores:} Seleccionados según el tipo de muestra y analito (ej. vidrio ámbar para fotosensibles, plástico para metales).
    \item \textbf{Condiciones de toma:} Temperatura, tiempo, agitación, etc.
    \item \textbf{Tamaño de muestra:} Suficiente para todos los análisis requeridos.
\end{itemize}


\resumebox{dato}{Control de calidad en alimentos}{\footnotesize
    El muestreo es crítico en la industria alimentaria para detectar contaminantes (como micotoxinas) o para verificar el contenido nutricional. Un muestreo inadecuado puede llevar a rechazar un lote bueno o, peor, a aprobar uno contaminado.
}